\section{Methodology}\label{sec:methodology}

The project adopts a design-science and User-Centered Design approach.
Design-science research focuses on building and evaluating artefacts that solve identified problems, while User-Centred Design (UCD) emphasises grounded design decisions in user needs and preferences.

BEACON was developed iteratively, by defining a set of ``Issues'', each with specific ``Acceptance Criteria''.
Each issue constitutes a feature or improvement to be implemented in the project.
For example, Alerts were implemented as a feature, with the acceptance criteria noting the need for finding a standard to prioritise and display alerts, and the need for a visual and audio component to the alert system.

\subsection{Research Paradigm}\label{subsec:research-paradigm}

This project is a qualitative study of supervisory interface design in the context of managing drone swarms in MSAR operations.

Qualitative analysis was used to influence how interface and swarm-behaviour decisions are designed and justified.
Choices about layout, alert hierarchy, swarm explainability, and interaction conventions were made by interpreting prior literature and applying it to BEACON's context.

Accordingly, the methodology focuses on explaining why the interface is structured in a particular way to support operator performance.
This ensures that theory-informed design reasoning is used to justify artefact decisions and evaluate their coherence with human-centred objectives (\cite{BABAEI2025103515}).

\subsection{Objectives and Hypotheses}\label{subsec:objectives-and-hypotheses}

To tie BEACON's design and evaluation to specific, verifiable hypotheses, three research questions were defined.

\begin{itemize}
    \item How effectively does BEACON's interface design improve situational awareness?
    \item To what extent does BEACON's interface design reduce operator workload, within acceptable limits for MSAR operations?
    \item How much do BEACON's interface optimisations to improve situational awareness improve mission success?
\end{itemize}

From these research questions, the following hypotheses can be derived:

\begin{itemize}
    \item BEACON's interface will provide users with adequate situational awareness, as inferred from performance metrics and self-reported scores.
    \item Workload, as measured by the National Aeronautics and Space Administration Task Load Index (NASA-TLX), will increase in high-stress scenarios, but will remain roughly within acceptable limits for MSAR operations.
    \item Missions supervised using BEACON's interface will have higher success rates, as measured by operational metrics like search coverage efficiency and target detection rate.
\end{itemize}

These research questions and hypotheses will provide the basis for the design and analysis of BEACON\@.

\subsection{HCI Approaches}\label{subsec:hci-approaches}

Several HCI frameworks and principles informed BEACON's design decisions.
Each principle was selected to address specific challenges arising from the supervisory-control context of drone-swarm management in an MSAR operation.

\subsubsection{Endsley's Situational Awareness Model}

Endsley's three-level model of Situational Awareness (SA) (\cite{situation_awareness}) was adopted as the primary theoretical lens for interface design.
Supervising an autonomous swarm across an open maritime area naturally requires operators to maintain awareness of multiple concurrent processes such as drone states and environmental conditions under time pressure.
Endsley's SA model provided a principled basis for deciding what information to surface and how that information should be evaluated.
For example, communication health is positioned as a persistent indicator near each drone, since losing track of it would represent a failure at SA's first level (perception).

\subsubsection{Proximity Compatibility Principle}

The Proximity Compatibility Principle (PCP)~\cite{doi:10.1518/001872095779049408} holds that information used together in a task should be displayed in close spatial proximity to each other to reduce the costs incurred from context-switching.
For example, in the context of BEACON, the central map is placed in proximity with the alert panel, so operators can cross-reference an event on the map with its associated alert without shifting focus across the screen.

\subsubsection{Aviation Crew Alerting Conventions}

Crew Alerting Systems structure alerts into tiers so that the most critical events are surfaced with priority without suppressing lower-priority context.
BEACON adopts this convention directly (\cite{RN3}) by implementing a four-tier hierarchy using aviation-standard colours.
This means alerts are grouped into red critical alerts, amber warnings, green advisories, and blue informational messages and are accompanied by category-specific audio cues.

This approach was chosen because studies like~\cite{CUMMINGS2007339} indicate that operators adopt load-shedding strategies under stress, ignoring streams of uniform alerts until a critical event is missed.
The hierarchy allows operators to triage issues quickly, while reducing the need for load-shedding.

\subsubsection{Explainability as a Design Constraint}

Research on supervisory control (\cite{AUTOMATIONEFFECTONUAVROUTING}) shows that overly opaque autonomy increases operator workload, because operators struggle to predict or analyse system behaviour.
BEACON treats explainability as a first-class design constraint.

For example, any swarm behaviour mechanism or path planning algorithm must be evaluated in terms of how easily its decisions can be explained to the operator.
This transparency reduces the need for micromanagement and supports higher-level supervisory engagement over reactive debugging.

\subsection{Development Approach, and Architecture}\label{subsec:development-approach-and-architecture}

BEACON was developed in React and TypeScript.
The map was rendered using a Canvas-based element, on which the drones, their paths, anomalies, and the mission boundary were drawn.

The architecture of the system emphasises domain purity, implementing clear separations between different, modular components of a larger system.
Logic like the environment generator, swarm behaviour, coverage planning, alerts, metrics collection, and more, were implemented as separate modules, hooking into a central \textit{MissionControl.tsx} context provider.

The main simulation loop runs in \textit{SimulationPage.tsx}, which uses the built-in \textit{requestAnimationFrame} function to update the state of the system at a fixed interval, to update drone kinematics, battery states, compute swarm adjustments, calculate per-drone communications, update sensor readings, and other related functions.
The interface is separated into different routes, like Landing, Mission Setup, Simulation, Results, NASA-TLX, and Management.
Transitions are handled with smooth animations, and the interface uses UI primitives like Badges, Fields, and ThemeToggles to provide consistent interaction patterns similar to other applications, while also providing a unique visual identity for BEACON\@.

\subsection{Practical and Technical Considerations}\label{subsec:practical-and-technical-considerations}

To keep the simulation responsive in the browser, and due to limitations in the development timeline, the physics and environmental effects are simplified.
This trade-off facilitates a focus on human-computer interaction (HCI) rather than physical modelling.

Related to the above, there is no real-time data link between physical drones and the simulation, which means that BEACON is best described as a Digital Model/Shadow prototype, rather than a full Digital Twin.
This decision was made due to the limited development and research timeline, with real-world physical testing being left for future work.
Human subject testing was also left for future work, due to ethical and logistical considerations, with the current evaluation framework relying on operational metrics and self-reported workload scores from the simulation itself.

Another important consideration was the choice of determinism versus variability in scenario generation.
To ensure that BEACON's evaluation is reproducible, and to allow for controlled comparisons between different interface designs, the environment generator works off of a pseudorandom seed-based generator, which allows for the same scenario to be generated across different runs of the simulation, through a common ``seed''.
Including coastal features was considered, but was ultimately left for future work to simplify environment generation and to focus on the core mechanics of drone supervision and interface design.

